% difstyle.tex -- red-pen markup for the dissertation diff (tools/dissertation-diff/build.py).
%
% build.py inserts this file at the end of the preamble of latexdiff's output,
% after latexdiff's own markup definitions, so everything here overrides them:
%
%   removed text      original ink with a red line through it
%   added text        red ink in a handwriting face (ECF Augie, from emerald)
%   new figure        red frame, "new figure" in the left margin
%   changed figure    old version shrunk and crossed out, new version framed
%   removed float     the original float, crossed out, caption struck through,
%                     with no number of its own so the edited column keeps the
%                     final numbering
%
% \DIFmark writes the page and height of every change to <jobname>.difpos,
% which build.py turns into the page's "next change" list. build.py's check
% builds (verify-old, verify-new) load <jobname>.mode at the end of this file to
% print only one version's text.

\usepackage{tikz}
\usepackage{emerald}
\definecolor{DIFpen}{RGB}{200,16,46}
% Removed text is printed in 1/255 grey: black to the eye, but verify.py can
% tell it apart from unchanged text when it reads the PDF back.
\definecolor{DIFdelink}{RGB}{1,1,1}
% Margin notes ("new figure", "removed") are one step off the pen colour for
% the same reason: they belong to neither version of the text.
\definecolor{DIFnote}{RGB}{200,16,47}

\makeatletter

% --- text ----------------------------------------------------------------
% Handwriting at 90% of the surrounding size so it sits on the same line.
\DeclareRobustCommand\DIFpenfont{%
  \fontfamily{fau}\fontseries{m}\fontshape{n}%
  \fontsize{\strip@pt\dimexpr 0.9\dimexpr\f@size pt\relax\relax}{\f@baselineskip}%
  \selectfont}
\newcommand\DIFstrikeon{%
  \bgroup\markoverwith{\textcolor{DIFpen}{\rule[0.52ex]{2pt}{0.8pt}}}\ULon}
\DeclareRobustCommand\DIFstrike[1]{%
  \ifmmode\text{\DIFstrikeon{\ensuremath{#1}}}\else\DIFstrikeon{#1}\fi}

\renewcommand{\DIFaddtex}[1]{{\protect\color{DIFpen}\ifmmode\else\protect\DIFpenfont\fi #1}}
\renewcommand{\DIFdeltex}[1]{{\protect\color{DIFdelink}\protect\DIFstrike{#1}}}
% build.py puts this between a removal and the addition that replaces it, so
% struck and handwritten words never touch.
\DeclareRobustCommand\DIFgap{\ifhmode\ifdim\lastskip=0pt\hskip0.3em\relax\fi\fi}

% --- change positions ----------------------------------------------------
\newwrite\DIFposout
\immediate\openout\DIFposout=\jobname.difpos\relax
\DeclareRobustCommand\DIFmark[1]{%
  \pdfsavepos
  \write\DIFposout{#1 \the\ReadonlyShipoutCounter\space\the\pdflastypos}}

\LetLtxMacro\DIFxaddbegin\DIFaddbegin
\LetLtxMacro\DIFxdelbegin\DIFdelbegin
\LetLtxMacro\DIFxaddbeginFL\DIFaddbeginFL
\LetLtxMacro\DIFxdelbeginFL\DIFdelbeginFL
\DeclareRobustCommand\DIFaddbegin{\DIFmark{a}\DIFxaddbegin}
\DeclareRobustCommand\DIFdelbegin{\DIFmark{d}\DIFxdelbegin}
\DeclareRobustCommand\DIFaddbeginFL{\DIFmark{a}\DIFxaddbeginFL}
\DeclareRobustCommand\DIFdelbeginFL{\DIFmark{d}\DIFxdelbeginFL}

% --- change IDs ------------------------------------------------------------
% build.py --overlay turns each "% REV <ID>" line into \DIFrevtag{<ID>}, which
% prints nothing: it writes the ID's page and position to <jobname>.difrev and
% the page draws the tag in the margin there. In the main vertical list (a
% marker between paragraphs) a whatsit could sit apart from the text it marks
% across a page break, so the tag waits for the next paragraph to start.
% Lists of contents and PDF bookmarks drop it.
\newwrite\DIFrevout
\immediate\openout\DIFrevout=\jobname.difrev\relax
\AtBeginDocument{\immediate\write\DIFrevout{@text
  \number\dimexpr 1in+\hoffset+\oddsidemargin\relax\space
  \number\dimexpr 1in+\hoffset+\oddsidemargin+\textwidth\relax}}
\newcommand\DIFrevwrite[1]{%
  \pdfsavepos
  \write\DIFrevout{#1 \the\ReadonlyShipoutCounter\space\the\pdflastxpos\space\the\pdflastypos}}
\gdef\DIFrevpending{}
\DeclareRobustCommand\DIFrevtag[1]{%
  \ifvmode\ifinner\DIFrevwrite{#1}%
    \else\g@addto@macro\DIFrevpending{\DIFrevwrite{#1}}\fi
  \else\DIFrevwrite{#1}\fi}
\AddToHook{para/begin}{%
  \ifx\DIFrevpending\@empty\else
    \DIFrevpending\global\let\DIFrevpending\@empty\fi}
\let\DIFxstarttoc\@starttoc
\def\@starttoc#1{\begingroup\let\DIFrevtag\@gobble\DIFxstarttoc{#1}\endgroup}
\pdfstringdefDisableCommands{\let\DIFrevtag\@gobble}

% --- figures -------------------------------------------------------------
\newsavebox\DIFfigbox
\newcommand\DIFmarginnote[2]{% #1 node, #2 note
  \node[overlay, anchor=north east, inner sep=0pt, text width=27mm, align=right,
        text=DIFnote, font=\DIFpenfont] at ([xshift=-9pt]#1.north west) {#2};}

% #1 note, #2 content: a red frame drawn outside the content (takes no space)
\newcommand\DIFframed[2]{%
  \sbox\DIFfigbox{#2}%
  \begin{tikzpicture}[baseline=(b.base)]
    \node[inner sep=0pt, outer sep=0pt] (b) {\usebox\DIFfigbox};
    \draw[overlay, DIFpen, line width=1pt, rounded corners=1.5pt]
      ([shift={(-3pt,-3pt)}]b.south west) rectangle ([shift={(3pt,3pt)}]b.north east);
    \DIFmarginnote{b}{#1}
  \end{tikzpicture}}

% #1 note, #2 content: framed and crossed corner to corner
\newcommand\DIFcrossed[2]{%
  \sbox\DIFfigbox{#2}%
  \begin{tikzpicture}[baseline=(b.base)]
    \node[inner sep=0pt, outer sep=0pt] (b) {\usebox\DIFfigbox};
    \draw[overlay, DIFpen, line width=0.8pt]
      ([shift={(-3pt,-3pt)}]b.south west) rectangle ([shift={(3pt,3pt)}]b.north east);
    \draw[overlay, DIFpen, line width=1.1pt, line cap=round]
      (b.south west) -- (b.north east) (b.north west) -- (b.south east);
    \DIFmarginnote{b}{#1}
  \end{tikzpicture}}

\renewcommand{\DIFaddincludegraphics}[2][]{%
  \DIFframed{new figure}{\DIFOincludegraphics[#1]{#2}}}
\renewcommand{\DIFdelincludegraphics}[2][]{%
  \DIFcrossed{removed}{\scalebox{0.6}{\DIFOincludegraphics[#1]{#2}}}}

% Same file name, new content: old version (from figures/difold/) shrunk and
% crossed out, then the new version framed.
\newcommand\DIFchangedgraphics[2][]{%
  \DIFmark{f}%
  \DIFcrossed{old version}{\DIFOincludegraphics[width=0.48\linewidth,
    height=0.26\textheight,keepaspectratio]{difold/#2}}\par
  \vspace{10pt}%
  \DIFframed{new version}{\DIFOincludegraphics[#1]{#2}}}

% --- removed floats --------------------------------------------------------
% #1 scale, #2 body (graphics or tabular); the margin note is \DIFcrossnote,
% which build.py sets inside a float that moved ("moved to C.1") or was
% rebuilt ("old version")
\newcommand\DIFcrossnote{removed}
\newcommand\DIFcrossbox[2]{%
  \DIFmark{f}%
  \DIFcrossed{\DIFcrossnote}{\color{DIFdelink}\scalebox{#1}{#2}}\par}
% a table whose rows or columns changed, in its new version (build.py's
% replace_restructured; the old one is a removed float just before it)
\newcommand\DIFnewtable[1]{\DIFmark{f}\DIFframed{new version}{#1}}
% a removed box that latexdiff could not mark (\figureplaceholder)
\newcommand\DIFdelbox[1]{\DIFmark{f}\DIFcrossed{removed}{#1}\par}
% #1 lof or lot, #2 figure or table, #3 old label text ("Table 3.3"),
% #4 old number, #5 caption, #6 list-of entry
\newcommand\DIFdelcaption[6]{%
  \caption*{\DIFdel{#3: #5}}%
  \addcontentsline{#1}{#2}{\protect\numberline{\DIFdel{#4}}\DIFdel{#6}}}
% the removed float's label: only the check build of the original sets it
\newcommand\DIFoldlabel[1]{}

% A float the markup makes taller than the page, scaled to fit (build.py's
% fit_floats); #1 is the scale.
\newsavebox\DIFfitbox
\newenvironment{DIFfit}[1]
  {\def\DIFfitscale{#1}\begin{lrbox}{\DIFfitbox}\begin{minipage}{\linewidth}}
  {\end{minipage}\end{lrbox}\centering\scalebox{\DIFfitscale}{\usebox\DIFfitbox}}

% Whole floats that were removed or added; build.py's check builds switch these off.
\newif\ifDIFshowremoved \DIFshowremovedtrue
\newif\ifDIFshowadded \DIFshowaddedtrue

\makeatother

% Check builds (verify-old, verify-new) bring their overrides in here.
\InputIfFileExists{\jobname.mode}{}{}
